\subsection*{MT-M30. An ideal-valued splitting of restriction to norm one}
Let \(K\) be a fixed number field, choose the norm splitting
and Haar normalization of MT-M24, and put
\[
 E_K=\mathcal S_{\mathrm E}(C_K),\qquad
 \mathcal I_a=|\cdot|^a\overline{\Theta_K}^{\,E_K},\qquad
 E_K^0=\{f\in E_K:f|_{C_K^1}=0\},
                                                               \tag{M30.1}
\]
where \(a\in\mathbf R\) is fixed. Let
\(\mathcal S(C_K^1)\) denote the Bruhat smooth functions on the
compact group \(C_K^1\). Restriction to norm one is the
continuous map
\[
                 r:E_K\longrightarrow\mathcal S(C_K^1).
 \]
There is a continuous linear, compact-translation-equivariant
map
\[
         B_a:\mathcal S(C_K^1)\longrightarrow\mathcal I_a,
                   \qquad rB_a=\operatorname{id}.             \tag{M30.2}
\]
No GRH input is needed. In particular
\[
 P_a=\operatorname{id}-B_ar
 \quad\hbox{is a continuous projection onto }E_K^0,\qquad
                   P_af-f\in\mathcal I_a.                    \tag{M30.3}
\]
Inclusion consequently induces a topological isomorphism
\[
        E_K^0/(\mathcal I_a\cap E_K^0)
                          \ \simeq\ E_K/\mathcal I_a,         \tag{M30.4}
\]
whose inverse is induced by \(P_a\).
The target in (M30.2) is the \emph{closed} centered
theta ideal. Actual summation-image membership of
an infinite character sum is not asserted.

Write
\[
         \mathcal D^0(C_K)
              =\{f\in\mathcal D(C_K):f|_{C_K^1}=0\}.
 \]
This space is dense in \(E_K^0\).
Therefore two continuous linear functionals on \(E_K\)
that annihilate \(\mathcal I_a\) are equal everywhere
if and only if they agree on \(\mathcal D^0(C_K)\).
The map \(P_a\) need not preserve compact radial
support or positivity.

\begin{proof}
\emph{Uniform normalized elements of the ideal.}
At an elementary compact level, use the family of
characters \(\beta_\nu\), \(\nu\in\mathbf Z^d\times F\),
and the parameter \(H_\nu=2+|\nu|\) from step~1 of
MT-M28. Each extension is trivial on the fixed norm
splitting. For \(\nu\ne0\) use the entire completion
\[
 \Xi_\nu(z)=Q_\nu^{z/2}
        \prod_{j=1}^{n_K}\Gamma_{\mathbf R}(z+\mu_{\nu,j})
                         L_\nu(z).
 \]
For \(\nu=0\) use the fixed pole-removed completion
of \(\zeta_K\), for example its standard completion
multiplied by \(z(z-1)\). These choices have precisely
the nontrivial completed zero divisors used by MT-M24.
All satisfy \(\Xi_\nu(2)\ne0\), by the Euler product
and the finite nonzero gamma factors there.
Define
\[
          F_\nu(z)=
              \frac{\Xi_\nu(z)}{\Xi_\nu(2)}\,e^{(z-2)^2}.
                                                               \tag{M30.5}
\]
It is entire, vanishes to the full order at every
completed zero, and satisfies \(F_\nu(2)=1\).
Let \(\psi_\nu\) be its inverse Mellin transform.
For each fixed \(\nu\), the fixed-strip bounds and
Gaussian decay in MT-M18 show
\(\psi_\nu\in\mathcal E\).

We need, and now prove, bounds uniform at this
compact level:
\[
           p_{N,j}(\psi_\nu)
                        \le C_{N,j}H_\nu^{B_{N,j}}.            \tag{M30.6}
\]
The finite conductor is uniformly bounded and
\(\operatorname{Re}\mu_{\nu,j}\) is a nonnegative
half-integer with \(|\mu_{\nu,j}|\ll H_\nu\).
The Euler lower bound at \(2\) and the uniform
finite-\(L\) polynomial upper bound are precisely
(M28.10) and (M28.12). Their proofs were
unconditional. On either line
\(\operatorname{Re}z=\pm(N+9/4)\), the same
gamma-ratio proof as in (M28.20), now with the
single fixed denominator point \(w=2\), gives
\[
 \left|\frac{\Xi_\nu(z)}{\Xi_\nu(2)}\right|
       \le C_N(H_\nu+|\operatorname{Im}z|)^{B_N}
                              e^{C_N|\operatorname{Im}z|}.
                                                               \tag{M30.7}
\]
In detail, shift numerator gamma factors right
by a fixed even integer. Every recurrence
denominator has real part a quarter-integer
of absolute value at least \(1/4\). The
remaining gamma ratio is in a fixed right
half-plane. The bounds
\(|\operatorname{Im}\psi(z)|\le C_c\) and
\(|\psi(z)|\le C_c+2\log(2+|z|)\) there
give (M30.7) by integrating the logarithmic
derivative in imaginary and real directions.
This uses a canceled ratio throughout, so
no factorial-size completion normalization
remains. The single principal character
has its fixed polynomial pole-removal
factor and is included by enlarging
the constants once.

Multiplication by the Gaussian in (M30.5)
turns (M30.7) into
\[
 |F_\nu(\pm(N+9/4)+it)|
                 \le C_NH_\nu^{B_N}e^{-c t^2}.
 \]
Inverting on these lines, differentiating
under the absolutely convergent integrals,
and using the positive line for \(x\ge0\)
and negative line for \(x\le0\), proves
(M30.6), just as in (M28.22)--(M28.23).
This part of MT-M28 used only Euler
convergence, the classical functional
equation and per-function finite order.
Neither its GRH contour selection nor
its narrow-strip minimum modulus was
used here.

Set
\[
                 \varphi_\nu(x)=e^{ax}\psi_\nu(x).
 \]
Multiplication by \(e^{ax}\) is a continuous
automorphism of \(\mathcal E\), so it
preserves the polynomial-in-\(H_\nu\)
form of (M30.6), with changed indices
and constants. Its normalization is
explicitly
\[
                   \mathcal M\varphi_\nu(2-a)=1.              \tag{M30.8}
\]
The full-zero-jet identification in
MT-M24 shows that this is an element
of the centered \(\nu\)-sector
theta ideal \(\mathcal I_{a,\nu}\).
One can also see the ideal membership
before taking its closure: the
Gaussian factor in (M30.5) is itself
the Mellin transform of an element
of \(\mathcal E\), so \(\psi_\nu\)
lies in the principal convolution
ideal generated by the inverse
transform of \(\Xi_\nu\).

\emph{A polynomial lower bound for the normalizing norm.}
It is important not to divide by an
uncontrolled small number. Put \(t_0=2-a\)
and choose an integer \(L>|t_0|+1\).
By (M30.6),
\[
                  p_{L,0}(\varphi_\nu)\le C H_\nu^B.
 \]
For \(R\ge1\), the tail of the integral
in (M30.8) satisfies
\[
 \left|\int_{|x|>R}e^{t_0x}\varphi_\nu(x)\,dx\right|
       \le C' H_\nu^B e^{-(L-|t_0|)R}.
 \]
Choose \(R=C_0+C_1\log H_\nu\), with
fixed constants large enough that
this bound is at most \(1/2\).
The integral over \([-R,R]\) has
absolute value at least \(1/2\).
Cauchy--Schwarz gives
\[
 \frac12\le
       (2R)^{1/2}e^{|t_0|R}\|\varphi_\nu\|_2 .
 \]
Since \(R\ll1+\log H_\nu\), it follows
that, for fixed \(c,C>0\),
\[
                   \|\varphi_\nu\|_2^2\ge cH_\nu^{-C}.        \tag{M30.9}
\]
This argument also covers \(t_0=0\).

For a radial function write
\(\varphi^*(x)=\overline{\varphi(-x)}\), and define
\[
              b_\nu
                =\frac{\varphi_\nu*\varphi_\nu^*}
                             {\|\varphi_\nu\|_2^2}.           \tag{M30.10}
\]
Its first convolution factor is in
\(\mathcal I_{a,\nu}\), and the second
is in \(\mathcal E\). Thus
\[
          b_\nu\in\mathcal I_{a,\nu},
          \qquad b_\nu(0)=1,\qquad
          p_{N,j}(b_\nu)\le C_{N,j}H_\nu^{C_{N,j}}.            \tag{M30.11}
\]
The value at zero is exactly the
squared \(L^2\) norm in (M30.10).
For the seminorm bound, use
\[
 p_{N,j}(f*g)
   \le p_{N,j}(f)\int_{\mathbf R}e^{N|y|}|g(y)|\,dy
   \le 2p_{N,j}(f)p_{N+1,0}(g)
 \]
and then (M30.6), (M30.9).
No stability of \(\mathcal I_a\)
under the involution for arbitrary
\(a\) is needed for this membership:
only its first factor and its
convolution-ideal property are used.

\emph{Summing over the compact character family.}
For \(g\in\mathcal S(C_K^1)\) at the
chosen compact level, write
\[
       g(c)=\sum_{\nu\in\Lambda}g_\nu\beta_\nu(c)^{-1},
       \qquad
       g_\nu=\int_{C_K^1}g(c)\beta_\nu(c)\,dc .
 \]
The coefficients decrease faster
than every power of \(H_\nu\).
Define, in the norm-splitting coordinates,
\[
 (B_ag)(c\,\tau(e^x))
       =\sum_{\nu\in\Lambda}
                   g_\nu\beta_\nu(c)^{-1}b_\nu(x).
                                                               \tag{M30.12}
\]
Each compact derivative costs a
fixed power of \(H_\nu\), and
(M30.11) bounds every weighted
radial derivative by a fixed power.
The rapid Fourier coefficients
therefore make (M30.12) absolutely
convergent in every seminorm of
the elementary \(E_K\) level.
The required source coefficient
bound uses finitely many compact
derivatives, so \(B_a\) is
continuous at that level.

The functions \(b_\nu\) were
defined from the character,
its fixed norm extension and
(M30.5), independently of any
choice of generators used to
prove the bounds. Thus the
definitions agree when levels
are enlarged. The elementary
inductive-limit definition of
the Bruhat smooth topology
now gives the asserted global
continuous map.
Every finite partial sum lies
in \(\mathcal I_a\), and this
ideal is closed in \(E_K\);
hence the whole sum lies in it.
At \(x=0\), (M30.11) gives
\(rB_ag=g\), with convergence
in the compact smooth topology.
The character formula also
shows that \(B_a\) commutes
with compact translations.
This proves (M30.2).

Equations (M30.3) follow from
\(rB_a=\operatorname{id}\).
If \(f\in E_K^0\), then \(P_af=f\).
If \(f\in\mathcal I_a\), then
\(P_af\in\mathcal I_a\cap E_K^0\).
The maps induced by inclusion
and \(P_a\) on the two quotient
spaces are therefore inverse.
Both are continuous by the
defining quotient topologies,
which proves (M30.4).

\emph{Compact test density and extension of equalities.}
For \(f\in E_K^0\), choose smooth
radial cutoffs \(\kappa_R(x)\)
equal to one on \([-R,R]\)
and zero outside \([-2R,2R]\).
The functions \(\kappa_R f\)
are in \(\mathcal D^0(C_K)\):
their compact-variable level
is unchanged, their norm
support is compact, and
\(\kappa_R(0)f(c)=0\).
They converge to \(f\) in
every weighted derivative
seminorm. Indeed the
differences and the derivatives
of the cutoffs are supported
where \(|x|\ge R\); a bound
at one larger exponential
weight gives an additional
factor tending to zero.
This proves density.

Suppose \(T_1,T_2\) are
continuous on \(E_K\), vanish
on \(\mathcal I_a\), and
agree on \(\mathcal D^0(C_K)\).
By density they agree on
\(E_K^0\). For every \(f\in E_K\),
\[
            T_i(f)=T_i(P_af),
 \]
since \(f-P_af\in\mathcal I_a\).
They are therefore equal
on all of \(E_K\). The
reverse implication is immediate.
\end{proof}

\paragraph{An exact restricted-trace criterion.}
The splitting gives a useful precise
test of the residual source claim.
For each spectral pair \((\chi,\rho)\)
choose real weights
\[
                   0\le w_{\chi,\rho}\le m_{\chi,\rho}.
 \]
Define
\[
 \begin{split}
 \tau_{\mathrm{full}}(f)
    &=\sum_{\chi,\rho}m_{\chi,\rho}
              \widehat f(\beta_\chi|\cdot|^{\,\rho-a}),\\
 \tau_w(f)
    &=\sum_{\chi,\rho}w_{\chi,\rho}
              \widehat f(\beta_\chi|\cdot|^{\,\rho-a}).
 \end{split}                                                    \tag{M30.13}
\]
These are continuous on \(E_K\) and
annihilate \(\mathcal I_a\),
unconditionally. To verify convergence,
at each compact level the unconditional
Jensen estimate (M28.13) counts
the nontrivial zeros in each unit
height window by
\(O(\log(H_\nu+|n|))\).
Their real parts are in the fixed
strip \([0,1]\). Integration by
parts in the compact variables
and in the norm variable bounds
their Mellin values by arbitrary
negative powers of \(H_\nu+|n|\),
uniformly on bounded sets at that
level. In the norm variable the
integrals of derivatives are
controlled by finitely many
exponential-Schwartz seminorms;
compact radial support is not
needed. The same summation as
in MT-M28 proves absolute
convergence and continuity.
Vanishing on \(\mathcal I_a\)
is the full zero-jet kernel
identity from MT-M24/MT-M27.

Consequently
\[
 \left.\tau_{\mathrm{full}}\right|_{\mathcal D^0(C_K)}
       =\left.\tau_w\right|_{\mathcal D^0(C_K)}
       \quad\Longleftrightarrow\quad
             w_{\chi,\rho}=m_{\chi,\rho}
                     \text{ for every }(\chi,\rho).           \tag{M30.14}
\]
For the nontrivial implication,
the splitting extends the equality
to \(E_K\). Fix one spectral
pair and use the finite-jet
interpolation of MT-M18 in
that single character sector
to get transform value one
at that zero and zero at
all other zeros, with every
unwanted higher jet zero.
Evaluating the difference of
the two functionals gives
\(m_{\chi,\rho}-w_{\chi,\rho}\).
It must vanish. The reverse
implication follows directly
from the sums.
Equivalently, any omitted
zero or lowered multiplicity
has a witness already in
\(\mathcal D^0(C_K)\):
apply \(P_a\) to this
finite-spectral interpolant,
then approximate by the
compact radial cutoffs above.
Continuity keeps the nonzero
difference for all sufficiently
large cutoffs.

In particular, at \(a=1/2\)
choose \(w_{\chi,\rho}=m_{\chi,\rho}\)
inside a specified open
centered strip and zero outside.
Equality of the full and
strip-restricted functionals
on tests vanishing at norm
one is equivalent to the
absence of every omitted
zero. Requiring this for
all positive strip widths
is exactly the Hecke
zero-location statement.
The test restriction does
not weaken that criterion.
This equivalence supplies
no proof of either side.

\paragraph{What this changes in the source comparison.}
Masullo \S1.4.6, original PDF p.~30,
restricts the proposed geometric
trace calculation to
\(\varphi(C_K^1)=0\);
\S1.7.2, pp.~34--35,
notes that a literal convolution
square \(h*h^*\) satisfying
this condition is zero.
That last assertion is correct:
at the identity
\[
                     (h*h^*)(1)=\|h\|_2^2.
 \]
The new splitting shows that
it is not an obstruction
at the corrected arithmetic
quotient level. The class
of \(h*h^*\) has the
representative \(P_{1/2}(h*h^*)\)
vanishing on \(C_K^1\), and
every continuous quotient
trace takes the same value
on these two representatives.
The new representative need
not itself be a convolution
square or have any pointwise
positivity property.

Let \(W_{\mathrm{raw},1/2}\)
be the classical raw Weil
distribution in the centering
convention of MT-M24. Its
spectral version is
\[
 W_{\mathrm{spec}}(f)
   =\widehat f(|\cdot|^{1/2})
      +\widehat f(|\cdot|^{-1/2})
      -W_{\mathrm{raw},1/2}(f)
   =\tau_{\mathrm{full}}(f).
\]
An exact classical trace
locator is Meyer,
\emph{On a representation of
the idele class group related
to primes and zeros of
\(L\)-functions},
Theorem~5.8, author pp.~43--44,
with the spectral identification
of Theorem~5.11, p.~47.
The complete supplied proof
of Theorem~5.8 and its
normalization were read in
this run; its upstream
trace-theoretic lemmas
remain cited classical inputs.
Here its continuation to
\(E_K\) can be defined by
the absolutely convergent
full spectral sum (M30.13),
so no unproved geometric
kernel continuity is inserted.

Theorem~12.5.1.1 in the
Masullo draft, original
pp.~519--520, asserts that
the same geometric expression
equals the trace with only
the zeros inside the chosen
strip, on the restricted
compact tests. These original
pages and \S1.4.6 were
inspected as page images
for MT-M29. By (M30.14),
the remaining gate is
exactly that geometric
identity independently
of zero location, rather
than admissibility of
positivity tests or a
missing extension from
the restricted test space.
The uniform ideal-valued
splitting is unconditional;
the value of the full
functional on convolution
squares is nonnegative
under \(A_D\) in the
Dirichlet family and
under \(M\) in the
general Hecke family,
as in MT-M19/MT-M28.
No positivity of the
corrected representatives
is inferred.

\paragraph{Run and framework locators.}
This is a new proof in
\texttt{GRH\_NIGHT/complexity},
\texttt{work/root/idele\_restriction\_splitting.tex}.
It advances the arithmetic
quotient comparison from
MT-M24/MT-M27/MT-M28 and
the determinant comparison
MT-M29. Its geometric
framework is the trace
functional on a quotient
by theta relations, beside
the inherited SA23
\texttt{geo-H-lidskii}
and \texttt{geo-H-genus}
rows, and cote~1
scans~79--80,111--112.
It supplies neither an
integral cycle comparison
nor an independent
Gan--Gross--Prasad
period positivity theorem.
